% !PRINTED 490
\weberpage{490}
2. \gesperrt{Ist}
\[
P = \pi,\ \pi',\ \pi'',\ \pi''' \ldots
\]
\gesperrt{eine Gruppe von Permutationen von $m$ Elementen und $\pi_1$ eine beliebige unter ihnen, so stimmt das System}
\[
P_1 = \pi\pi_1,\ \pi'\pi_1,\ \pi''\pi_1,\ \pi'''\pi_1 \ldots,
\]
\gesperrt{von der Reihenfolge abgesehen, vollständig mit $P$ überein.}

Denn die in $P_1$ enthaltenen Permutationen, deren Anzahl ebenso gross ist, wie die der Permutationen $P$, sind wegen der Gruppeneigenschaft von $P$ (\S.\,148, 3.) jedenfalls alle unter den Permutationen von $P$ enthalten. Ausserdem sind die Permutationen von $P_1$ (nach \S.\,148, 2.) alle von einander verschieden, und also ist $P$ mit $P_1$, von der Reihenfolge abgesehen, identisch. Ebenso kann man auch zeigen, dass
\[
P_1' = \pi_1\pi,\ \pi_1\pi',\ \pi_1\pi'',\ \pi_1\pi''', \ldots
\]
mit $P$ identisch ist.

Hieraus ergeben sich die folgenden Sätze:

3. \gesperrt{Jede Permutationsgruppe enthält die identische Permutation.}

Denn unter den Permutationen von $P_1$ muss auch $\pi_1$ selbst vorkommen und wenn $\pi_0\pi_1 = \pi_1$ ist, so ist (nach \S.\,148, 2.) $\pi_0$ die identische Permutation.

4. \gesperrt{Eine Permutationsgruppe $P$ enthält zu jeder Permutation auch die entgegengesetzte.}

Denn in $P_1$ muss nach 3. die identische Permutation vorkommen. Ist aber $\pi\pi_1$ die identische Permutation, so ist $\pi = \pi_1^{-1}$. Hierin kann $\pi_1$ jede Permutation aus $P$ sein.

Wenn an Stelle der Variablen $u, u_1 \ldots u_{m-1}$ bestimmte Grössen $\alpha, \alpha_1, \alpha_2 \ldots \alpha_{m-1}$, seien es Zahlenwerthe oder Grössen irgend eines Körpers, gesetzt werden, so verliert der Satz 1. seine allgemeine Gültigkeit. Nehmen wir z.\,B. $m = 3$ und setzen zwischen den drei Grössen $\alpha, \alpha_1, \alpha_2$ die Relation $2\alpha_2 = \alpha + \alpha_1$
fest, indem wir sie sonst nicht weiter beschränken, so bleibt die Function $\alpha_1 - \alpha_2$ ungeändert nur durch die identische Permutation und durch die Permutation
\[
\begin{pmatrix} 0, & 1, & 2 \\ 1, & 2, & 0 \end{pmatrix}.
\]
